% GENERATED FILE. Edit canonical lessons, then run npm run generate:books.
% canonical-source-hash: fnv1a64:a5a2d3423ffd0ddd
% canonical-lessons: TA-C448-cemi, TA-C448-celavali, TA-C448-vinakku, TA-C448-mincu, TA-C448-varumanam

\chapter{Saving and Spending}
\label{ch:ta-saving-and-spending}
\hlchaptermodality{448}

\begin{chapteropening}
\textbf{By the end of this chapter:} \emph{I can talk about saving, spending and wasting money, what is left over, and income.}

Complete the last lesson of chapter 448: I can talk about saving, spending and wasting money, what is left over, and income.
\end{chapteropening}

\section[\texorpdfstring{\ta{சேமி} (cēmi)}{சேமி (cēmi)}]{\ta{சேமி} (cēmi) — to save up; to store}

\label{lesson:TA-C448-cemi}



\begin{quote}
Before the new one: say the Tamil for the remainder, then the Tamil for a sum.
\end{quote}

\subsection*{You'll want to know: \ta{சேமி}}
\textbf{\ta{சேமி}} — \emph{cēmi} — \textquotedblleft{}to save up; to store\textquotedblright{}.

\begin{quote}
\ta{நான்} \ta{பயணத்துக்காகப்} \ta{பணம்} \ta{சேமித்தேன்}.
\end{quote}

\begin{quote}
\emph{nāṉ payaṇattukkākap paṇam cēmittēṉ.}
\end{quote}

\begin{quote}
I saved money for the trip.
\end{quote}

Past (I …): \ta{சேமித்தேன்}, \emph{cēmittēṉ}.

\ta{சேமிப்பு}, savings, is its noun.

\begin{cousinweb}
\textbf{\ta{சேமி}} (\emph{cēmi}) is usually traced to Sanskrit \emph{kṣema}, \textquotedblleft{}safety, welfare\textquotedblright{} (CDIAL 3745), through Tamil \ta{சேமம்}, \textquotedblleft{}safekeeping\textquotedblright{}.
\end{cousinweb}

\begin{grammarlens}[title={What you've built}]
One more word for getting around.
\end{grammarlens}

\subsection*{Your turn}
\begin{itemize}
\raggedright
  \item \emph{Say it:} \emph{cēmi}
  \item \emph{Say it:} \emph{cēmi}, once more
  \item \emph{Say it:} \emph{nāṉ payaṇattukkākap paṇam cēmittēṉ.} — I saved money for the trip.
  \item \emph{Recall:} say the Tamil for the remainder, then the Tamil for a sum, then say \emph{cēmi} again
\end{itemize}

\begin{tcolorbox}[breakable,colback=teal!4,colframe=teal!35!black,title={Before you move on}]
What does \ta{சேமி} mean? (\textquotedblleft{}To save up; to store\textquotedblright{}.) Say it once more.
\end{tcolorbox}
\section[\texorpdfstring{\ta{செலவழி} (celavaḻi)}{செலவழி (celavaḻi)}]{\ta{செலவழி} (celavaḻi) — to spend money; to spend}

\label{lesson:TA-C448-celavali}



\begin{quote}
Before the new one: say the Tamil for a sum, then the Tamil for to save up.
\end{quote}

\subsection*{You'll want to know: \ta{செலவழி}}
\textbf{\ta{செலவழி}} — \emph{celavaḻi} — \textquotedblleft{}to spend money; to spend\textquotedblright{}.

\begin{quote}
\ta{விடுமுறையில்} \ta{நாங்கள்} \ta{நிறைய} \ta{பணம்} \ta{செலவழித்தோம்}.
\end{quote}

\begin{quote}
\emph{viṭumuṟaiyil nāṅkaḷ niṟaiya paṇam celavaḻittōm.}
\end{quote}

\begin{quote}
We spent a lot of money on holiday.
\end{quote}

Past (I …): \ta{செலவழித்தேன்}, \emph{celavaḻittēṉ}.

\ta{செலவு}, an expense, + \ta{அழி}, to use up.

\begin{cousinweb}
\textbf{\ta{செலவழி}} (\emph{celavaḻi}) is built from \ta{செலவு}, \textquotedblleft{}going, expenditure\textquotedblright{} (DEDR 2781), and \ta{அழி}, \textquotedblleft{}to perish, to be used up\textquotedblright{} (DEDR 277).
\end{cousinweb}

\begin{grammarlens}[title={What you've built}]
One more word for getting around.
\end{grammarlens}

\subsection*{Your turn}
\begin{itemize}
\raggedright
  \item \emph{Say it:} \emph{celavaḻi}
  \item \emph{Say it:} \emph{celavaḻi}, once more
  \item \emph{Say it:} \emph{viṭumuṟaiyil nāṅkaḷ niṟaiya paṇam celavaḻittōm.} — We spent a lot of money on holiday.
  \item \emph{Recall:} say the Tamil for a sum, then the Tamil for to save up, then say \emph{celavaḻi} again
\end{itemize}

\begin{tcolorbox}[breakable,colback=teal!4,colframe=teal!35!black,title={Before you move on}]
What does \ta{செலவழி} mean? (\textquotedblleft{}To spend money; to spend\textquotedblright{}.) Say it once more.
\end{tcolorbox}
\section[\texorpdfstring{\ta{வீணாக்கு} (vīṇākku)}{வீணாக்கு (vīṇākku)}]{\ta{வீணாக்கு} (vīṇākku) — to waste}

\label{lesson:TA-C448-vinakku}



\begin{quote}
Before the new one: say the Tamil for to save up, then the Tamil for to spend money.
\end{quote}

\subsection*{You'll want to know: \ta{வீணாக்கு}}
\textbf{\ta{வீணாக்கு}} — \emph{vīṇākku} — \textquotedblleft{}to waste\textquotedblright{}.

\begin{quote}
\ta{தண்ணீரை} \ta{வீணாக்கக்} \ta{கூடாது}.
\end{quote}

\begin{quote}
\emph{taṇṇīrai vīṇākkak kūṭātu.}
\end{quote}

\begin{quote}
You must not waste water.
\end{quote}

Past (I …): \ta{வீணாக்கினேன்}, \emph{vīṇākkiṉēṉ}.

\ta{வீண்}, useless, + \ta{ஆக்கு}, to make.

\begin{cousinweb}
\textbf{\ta{வீணாக்கு}} (\emph{vīṇākku}) is built from \ta{வீண்}, \textquotedblleft{}uselessness\textquotedblright{} (DEDR 5454), and \ta{ஆக்கு}, \textquotedblleft{}to make\textquotedblright{} (DEDR 333).
\end{cousinweb}

\begin{grammarlens}[title={What you've built}]
One more word for getting around.
\end{grammarlens}

\subsection*{Your turn}
\begin{itemize}
\raggedright
  \item \emph{Say it:} \emph{vīṇākku}
  \item \emph{Say it:} \emph{vīṇākku}, once more
  \item \emph{Say it:} \emph{taṇṇīrai vīṇākkak kūṭātu.} — You must not waste water.
  \item \emph{Recall:} say the Tamil for to save up, then the Tamil for to spend money, then say \emph{vīṇākku} again
\end{itemize}

\begin{tcolorbox}[breakable,colback=teal!4,colframe=teal!35!black,title={Before you move on}]
What does \ta{வீணாக்கு} mean? (\textquotedblleft{}To waste\textquotedblright{}.) Say it once more.
\end{tcolorbox}
\section[\texorpdfstring{\ta{மிஞ்சு} (miñcu)}{மிஞ்சு (miñcu)}]{\ta{மிஞ்சு} (miñcu) — to be left over}

\label{lesson:TA-C448-mincu}



\begin{quote}
Before the new one: say the Tamil for to spend money, then the Tamil for to waste.
\end{quote}

\subsection*{You'll want to know: \ta{மிஞ்சு}}
\textbf{\ta{மிஞ்சு}} — \emph{miñcu} — \textquotedblleft{}to be left over\textquotedblright{}.

\begin{quote}
\ta{பயணம்} \ta{முடிந்தபோது} \ta{கொஞ்சம்} \ta{பணம்} \ta{மிஞ்சியது}.
\end{quote}

\begin{quote}
\emph{payaṇam muṭintapōtu koñcam paṇam miñciyatu.}
\end{quote}

\begin{quote}
When the trip was over, a little money was left over.
\end{quote}

Past (it …): \ta{மிஞ்சியது}, \emph{miñciyatu}.

Also to outdo: \ta{அவரை} \ta{யாரும்} \ta{மிஞ்ச} \ta{முடியாது}, no one can outdo him.

\begin{cousinweb}
\textbf{\ta{மிஞ்சு}} (\emph{miñcu}) is native Dravidian (DEDR 4838), \textquotedblleft{}to exceed, to increase, to remain over\textquotedblright{}.
\end{cousinweb}

\begin{grammarlens}[title={What you've built}]
One more word for getting around.
\end{grammarlens}

\subsection*{Your turn}
\begin{itemize}
\raggedright
  \item \emph{Say it:} \emph{miñcu}
  \item \emph{Say it:} \emph{miñcu}, once more
  \item \emph{Say it:} \emph{payaṇam muṭintapōtu koñcam paṇam miñciyatu.} — When the trip was over, a little money was left over.
  \item \emph{Recall:} say the Tamil for to spend money, then the Tamil for to waste, then say \emph{miñcu} again
\end{itemize}

\begin{tcolorbox}[breakable,colback=teal!4,colframe=teal!35!black,title={Before you move on}]
What does \ta{மிஞ்சு} mean? (\textquotedblleft{}To be left over\textquotedblright{}.) Say it once more.
\end{tcolorbox}
\section[\texorpdfstring{\ta{வருமானம்} (varumāṉam)}{வருமானம் (varumāṉam)}]{\ta{வருமானம்} (varumāṉam) — income; earnings}

\label{lesson:TA-C448-varumanam}



\begin{quote}
Before the new one: say the Tamil for to waste, then the Tamil for to be left over.
\end{quote}

\subsection*{You'll want to know: \ta{வருமானம்}}
\textbf{\ta{வருமானம்}} — \emph{varumāṉam} — \textquotedblleft{}income; earnings\textquotedblright{}.

\begin{quote}
\ta{அவருடைய} \ta{வருமானம்} \ta{குறைவு}, \ta{ஆனால்} \ta{செலவு} \ta{அதிகம்}.
\end{quote}

\begin{quote}
\emph{avaruṭaiya varumāṉam kuṟaivu, āṉāl celavu atikam.}
\end{quote}

\begin{quote}
His income is low, but his expenses are high.
\end{quote}

\ta{மாத} \ta{வருமானம்} is a monthly income.

\begin{cousinweb}
\textbf{\ta{வருமானம்}} (\emph{varumāṉam}) is built on \emph{varu-}, the stem of \ta{வா}, \textquotedblleft{}to come\textquotedblright{} (DEDR 5270), like \ta{வருகை}; the history of its ending is not settled.
\end{cousinweb}

\begin{grammarlens}[title={What you've built}]
One more word for getting around.
\end{grammarlens}

\subsection*{Your turn}
\begin{itemize}
\raggedright
  \item \emph{Say it:} \emph{varumāṉam}
  \item \emph{Say it:} \emph{varumāṉam}, once more
  \item \emph{Say it:} \emph{avaruṭaiya varumāṉam kuṟaivu, āṉāl celavu atikam.} — His income is low, but his expenses are high.
  \item \emph{Recall:} say the Tamil for to waste, then the Tamil for to be left over, then say \emph{varumāṉam} again
\end{itemize}

\begin{tcolorbox}[breakable,colback=teal!4,colframe=teal!35!black,title={Before you move on}]
What does \ta{வருமானம்} mean? (\textquotedblleft{}Income; earnings\textquotedblright{}.) Say it once more.
\end{tcolorbox}
